\hmsubsection{Hardware Architecture}{UMA}{}
The hardware architecture of the system is organised into three 
functional layers. The perception layer consists of the OAK-D S2 
camera, which captures visual data of the workspace. The decision 
layer consists of the Raspberry Pi~5, which executes the image 
processing pipeline and the inverse kinematics solver. The actuation 
layer consists of the OpenRB-150 microcontroller and five Dynamixel 
servos, which translate motor commands into physical movement of 
the arm. A fourth, transverse subsystem provides regulated electrical 
power to all components.

This subsection describes each part of the hardware architecture in 
turn. Section~7.2.1 describes the components of the vision system. 
Sections~7.2.2--7.2.5 describe the power architecture, including 
the theoretical current demand, the selection of the power supply, 
and the topology used to distribute current to the motors. 
Sections~7.2.6--7.2.9 describe the selection and sizing of the 
Dynamixel servos for each joint, and Sections~7.2.10--7.2.13 
describe the selection of the OpenRB-150 as the embedded motor 
controller.



\hmsubsubsection{Hardware Components for the Vision System}{OAH}{}

This subsection presents the hardware components that constitute the vision system used in
this project. These components provide the necessary functionality for image acquisition,
processing, and data transfer, enabling reliable detection and classification of objects based on
colour.

\begin{table}[htbp]
    \centering
    \renewcommand{\arraystretch}{1.3}
    \begin{tabular}{|p{2.3in}|p{2.1in}|p{2.1in}|}
        \hline
        \textbf{Component} & \textbf{Description} & \textbf{Purpose} \\
        \hline
        OAK-D S2 Camera (Luxonis) & 
        Advanced AI camera with integrated depth sensing and onboard processing capabilities &
        Captures image data and performs real-time object detection of red and blue balls \\
        \hline
        Raspberry Pi 5 & 
        Single-board computer with sufficient processing power for embedded applications &
        Executes the image processing pipeline and manages communication with other system components \\
        \hline
        microSD Card (128GB) & 
        Storage device used for the operating system and software &
        Stores the operating system, required dependencies, and project code necessary for execution \\
        \hline
        USB 3 Cable (USB-A to USB-C) & 
        High-speed data transfer cable &
        Connects the camera to the Raspberry Pi and ensures sufficient bandwidth for real-time data transmission \\
        \hline
    \end{tabular}
    \caption{Hardware components of the vision system}
    \label{tab:vision_hardware}
\end{table}

The components listed in Table~\ref{tab:vision_hardware} constitute the complete vision system. The OAK-D S2 camera performs image acquisition and initial onboard processing, while the Raspberry Pi 5 serves as the primary processing unit responsible for executing the image pipeline and coordinating communication with the rest of the system.

The microSD card provides local storage for the operating system and required software, enabling independent operation of the platform. In addition, the USB 3 connection ensures sufficient bandwidth for real-time transfer of image data between the camera and the processing unit.

Overall, this hardware configuration provides a modular and scalable foundation for the vision system, supporting both the current implementation and future improvements in processing performance and detection accuracy.



\hmsubsubsection{Power Architecture}{UMA}{}

The arm uses five Dynamixel servos that operate from a 12~V supply. 
The selection and arrangement of the power supply was identified 
early in the project as a critical design concern, since the previous 
iteration of the project encountered repeated power-related failures 
and required unstable workarounds during testing \cite{hivemind}. The 
present design therefore aims to provide a stable supply with 
sufficient current capacity, and to distribute current safely across 
the motor chain.

\hmsubsubsection{Theoretical Current Demand}{UMA}{}
Each Dynamixel servo has a published stall current that represents 
the maximum current it draws when the motor is mechanically blocked. 
While stall conditions are not expected during normal operation, the 
total stall current across all five servos provides an upper bound 
on the worst-case demand. The values given in the manufacturer's 
datasheets \cite{robotis_dynamixel} are summarised in 
Table~\ref{tab:stall_currents}. The choice of a Dynamixel servo for 
the claw, rather than the analogue 5~V servo used in the previous 
iteration, is discussed in Section~7.2.7.

\begin{table}[htbp]
    \centering
    \renewcommand{\arraystretch}{1.3}
    \begin{tabular}{|l|l|c|}
        \hline
        \textbf{Joint} & \textbf{Servo} & \textbf{Stall Current (A)} \\
        \hline
        Base (M1)     & XM430-W350 & 2.3 \\
        \hline
        Shoulder (M2) & XM540-W270 & 4.4 \\
        \hline
        Elbow (M3)    & XM430-W350 & 2.3 \\
        \hline
        Wrist (M4)    & XL430-W250 & 1.3 \\
        \hline
        Claw (M5)     & XM430-W350 & 2.3 \\
        \hline
        \multicolumn{2}{|r|}{\textbf{Total}} & \textbf{12.6} \\
        \hline
    \end{tabular}
    \caption{Stall current per Dynamixel servo. The total represents 
    the theoretical worst case in which all five motors stall 
    simultaneously.}
    \label{tab:stall_currents}
\end{table}

The total worst-case demand of 12.6~A is rarely encountered in 
practice. During typical pick-and-place operation, the arm draws 
approximately 2--5~A, since only a subset of motors are accelerating 
at any given time and the payload is light (a 50~mm ball weighing 
approximately 20~g). The stall current is relevant primarily for 
sizing the supply to tolerate transient peaks and short-duration 
overloads.

\hmsubsubsection{Selection of the Power Supply}{UMA}{}
A dedicated 12~V battery rated at 10~A was initially considered, as 
it would provide full coverage of the worst-case current demand and 
allow the system to operate independently of mains power, which is 
relevant for the future integration of the arm onto a mobile platform. 
However, after consultation with the project supervisor, it was 
decided to retain the laboratory bench power supply rated at 
12~V/6~A as the final solution. This decision was based on three 
considerations. First, the laboratory power supply is regulated and 
protected against over-current and short circuits, which simplifies 
the electrical design. Second, the typical operating current of 
2--5~A is well within the 6~A capacity of the supply. Third, the 
mobile-platform integration is outside the scope of the current 
project, and a battery solution can be introduced in a later 
iteration without affecting the rest of the architecture.

\hmsubsubsection{Distribution Topology}{UMA}{}
A single power supply must be distributed to five servos that are 
physically arranged along the arm. The previous iteration of the 
project distributed power through a single daisy-chained TTL bus, 
in which all current was passed through the JST connector of each 
servo before reaching the next \cite{hivemind}. This arrangement was 
found to be problematic for two reasons. The first servo in the 
chain carries the cumulative current of all downstream motors. For 
the worst-case demand of approximately 12~A, this exceeds the rated 
capacity of standard 3-pin JST connectors, which are typically 
specified for 3--5~A. The second concern is voltage drop along the 
chain: each connector and cable segment introduces a small resistive 
loss, which accumulates and reduces the supply voltage available to 
the servos at the end of the chain.

To address these issues, the present design uses a two-branch star 
topology with two separate cables leaving the OpenRB-150 controller. 
The first branch carries the base (M1) and shoulder (M2) servos, 
which together account for the largest share of the total current 
demand (6.7~A combined stall). The second branch carries the elbow 
(M3), wrist (M4), and claw (M5) servos, with a combined stall 
current of 5.9~A. Within each branch, the servos remain 
daisy-chained, since the cumulative current within either branch is 
within the rated capacity of the JST connector. The data line 
follows the same physical wiring, as the Dynamixel TTL protocol 
shares the connector pins with the power rails.

The resulting topology is shown in 
Figure~\ref{fig:power_topology}. This arrangement preserves the 
convenience of the daisy-chain protocol for command transmission 
while halving the current load on the first connector of either 
branch. It also reduces voltage drop at the servos furthest from 
the controller, since each branch is shorter than a single 
five-servo chain.

\begin{figure}[H]
    \centering
    \begin{tikzpicture}[
        node distance=0.6cm and 1.0cm,
        psu/.style={
            rectangle, draw, thick, 
            minimum width=2.4cm, minimum height=0.9cm,
            fill=gray!15, font=\small\bfseries
        },
        controller/.style={
            rectangle, draw, thick, rounded corners=2pt,
            minimum width=2.4cm, minimum height=0.9cm,
            fill=blue!10, font=\small\bfseries
        },
        motor/.style={
            rectangle, draw, thick, rounded corners=2pt,
            minimum width=1.8cm, minimum height=1.0cm,
            fill=orange!15, align=center, font=\small
        },
        powerline/.style={
            -{Latex[length=2mm]}, very thick, red!70!black
        },
        chainline/.style={
            -{Latex[length=2mm]}, thick, red!70!black, dashed
        },
        branchlabel/.style={
            font=\scriptsize\itshape, text=black!70
        }
    ]

    % Power supply
    \node[psu] (psu) {12 V / 6 A PSU};
    
    % OpenRB-150 controller below PSU
    \node[controller, below=of psu] (openrb) {OpenRB-150};
    
    % Branch A: M1 -> M2 (top branch)
    \node[motor, below left=1.2cm and 0.8cm of openrb] (m1) {M1\\Base\\\scriptsize XM430};
    \node[motor, right=of m1] (m2) {M2\\Shoulder\\\scriptsize XM540};
    
    % Branch B: M3 -> M4 -> M5 (bottom branch)
    \node[motor, below=2.2cm of m1] (m3) {M3\\Elbow\\\scriptsize XM430};
    \node[motor, right=of m3] (m4) {M4\\Wrist\\\scriptsize XL430};
    \node[motor, right=of m4] (m5) {M5\\Claw\\\scriptsize XM430};
    
    % PSU -> OpenRB
    \draw[powerline] (psu) -- (openrb);
    
    % Branch A: OpenRB -> M1 -> M2
    \draw[powerline] (openrb.south) 
        -- ++(0,-0.4) 
        -| (m1.north);
    \draw[chainline] (m1) -- (m2);
    
    % Branch B: OpenRB -> M3 -> M4 -> M5
    \draw[powerline] (openrb.south) 
        -- ++(0,-0.4) 
        -| (m3.north);
    \draw[chainline] (m3) -- (m4);
    \draw[chainline] (m4) -- (m5);
    
    % Branch labels
    \node[branchlabel, right=0.1cm of m2] {Branch A: 6.7 A};
    \node[branchlabel, right=0.1cm of m5] {Branch B: 5.9 A};
    
    % Legend
    \begin{scope}[shift={(-3.8,-5.8)}]
        \draw[powerline] (0,0) -- (0.7,0) node[right, font=\scriptsize] 
            {Dedicated cable from controller};
        \draw[chainline] (0,-0.4) -- (0.7,-0.4) node[right, font=\scriptsize] 
            {Daisy-chained connection};
    \end{scope}

    \end{tikzpicture}
    \caption[Two-branch star topology for distributing 12\,V power across five Dynamixel servos.]{Two-branch star topology used to distribute 12~V power 
    from the laboratory supply across the five Dynamixel servos. Each 
    branch leaves the OpenRB-150 through a separate cable, halving 
    the peak current load on the first connector of either chain. The 
    branch currents shown are the worst-case stall totals for each 
    branch.}
    \label{fig:power_topology}
\end{figure}

\hmsubsubsection{Motor Selection and Joint Sizing}{UMA}{}

The selection of actuators determines both the mechanical 
capabilities of the arm and the complexity of the control system. 
Five servo motors were required, one for each degree of freedom 
plus the gripper. The selection process therefore involved two 
decisions: the choice of a motor family suitable for the 
application, and the sizing of individual motors for each joint 
according to the load they are expected to carry.

\hmsubsubsection{Choice of Motor Family}{UMA}{}
The Dynamixel X-series from ROBOTIS was selected as the motor 
family for all five joints. Three properties of this family were 
particularly relevant for the project. First, each motor 
incorporates an integrated controller that provides closed-loop 
position control with feedback from a built-in encoder, eliminating 
the need for external motor drivers and feedback circuitry. Second, 
all motors in the family communicate over a shared TTL bus using 
the Dynamixel Protocol 2.0, allowing up to 253 motors to be 
addressed individually through a single daisy-chained connection 
\cite{robotis_dynamixel}. Third, the family includes models with 
different torque ratings that share the same communication 
protocol, mounting interface, and software API, enabling 
heterogeneous motor selection across joints without increasing 
software complexity.

The previous iteration of the project used a Dynamixel-based arm 
for joints~1--4 but employed a separate analogue 5~V servo for 
the gripper, controlled through a PCA9685 PWM driver connected to 
an Arduino Mega \cite{hivemind}. This arrangement required two 
distinct control paths, two separate power supplies, and 
additional driver circuitry. In the present design, the analogue 
gripper servo was replaced with a Dynamixel XM430-W350. This 
substitution eliminates the PCA9685 driver, the second power 
supply, and the dedicated PWM control path, allowing all five 
motors to be driven through the same TTL bus and the same firmware 
routine. The simplification improves system reliability and 
reduces the number of failure points in the electrical design.

\hmsubsubsection{Joint-Specific Motor Sizing}{UMA}{}
The torque required at each joint depends on the mass of the 
downstream links and the horizontal distance from the joint axis 
to the centre of mass. For a serial arm, joints closer to the 
base must support a longer lever arm and therefore experience 
higher torque. The motors were sized accordingly, with stronger 
models placed at the joints subject to the greatest mechanical 
load. The selected models and their key specifications are 
summarised in Table~\ref{tab:motor_selection}.

\begin{table}[htbp]
    \centering
    \renewcommand{\arraystretch}{1.3}
    \begin{tabular}{|l|l|c|c|c|}
        \hline
        \textbf{Joint} & \textbf{Model} & 
        \textbf{Stall Torque} & 
        \textbf{Mass} &
        \textbf{Cost} \\
         & & \textbf{(N$\cdot$m)} & \textbf{(g)} & \textbf{(rel.)} \\
        \hline
        Base (M1)     & XM430-W350 & 4.1 & 82  & Medium \\
        \hline
        Shoulder (M2) & XM540-W270 & 10.6 & 165 & High \\
        \hline
        Elbow (M3)    & XM430-W350 & 4.1 & 82  & Medium \\
        \hline
        Wrist (M4)    & XL430-W250 & 1.5 & 57  & Low \\
        \hline
        Claw (M5)     & XM430-W350 & 4.1 & 82  & Medium \\
        \hline
    \end{tabular}
    \caption{Selected Dynamixel servos for each joint, with key 
    specifications from the manufacturer datasheets 
    \cite{robotis_dynamixel}.}
    \label{tab:motor_selection}
\end{table}

The shoulder joint (M2) supports the entire mass of the forearm, 
wrist, gripper, and any payload, multiplied by the full horizontal 
reach of the arm. It therefore experiences the highest torque of 
all five joints. The XM540-W270, with a stall torque of 
10.6~N$\cdot$m, was selected for this joint to provide sufficient 
margin above the worst-case loaded condition.

The elbow joint (M3) supports the mass of the wrist, gripper, and 
payload over the length of the forearm. The XM430-W350 was chosen 
as a compromise between torque capacity (4.1~N$\cdot$m) and weight 
(82~g), since the elbow motor itself contributes to the load on 
the shoulder. Selecting a heavier motor for the elbow would have 
required upgrading the shoulder motor as well.

The base joint (M1) does not lift the arm but rotates it 
horizontally, and therefore experiences torque only from inertial 
acceleration during base motion. The XM430-W350 was chosen for 
consistency with the elbow, since both joints share the same 
mounting and connector interface.

The wrist joint (M4) carries only the gripper and payload, and 
operates close to its rotation axis. The XL430-W250, with a stall 
torque of 1.5~N$\cdot$m, was selected for this joint. Its lower 
mass (57~g) and lower cost compared with the XM430 reduce both 
the load on the elbow motor and the overall cost of the arm 
without compromising functionality.

The gripper (M5) requires sufficient torque to close around a 
50~mm ball with a firm grip but does not carry the load of any 
downstream joints. The XM430-W350 was selected for the gripper, 
matching the base and elbow motors and allowing all four 
non-shoulder joints to share the same spare-part inventory.

\hmsubsubsection{Link Lengths}{UMA}{}
The arm's link lengths were measured directly from the assembled 
mechanical structure as $L_1 = 25.5$~cm (shoulder to elbow), 
$L_2 = 23.0$~cm (elbow to wrist), and $L_3 = 22.0$~cm (wrist to 
claw tip). The first two links determine the planar reach of the 
arm, giving a theoretical maximum of $L_1 + L_2 = 48.5$~cm. The 
third link contributes to the vertical position of the claw tip 
but does not extend the planar reach, since the wrist always 
orients the gripper to point directly downwards. The geometric 
basis for this configuration is described in 
Section~8.7.

The combined link length of approximately 70.5~cm is longer than 
the 52--54~cm initially planned for the arm. The increase is 
primarily due to the third link ($L_3$), which was lengthened 
during mechanical design to accommodate the wrist-mounted OAK-D~S2 
camera and the gripper assembly. The longer $L_3$ does not affect 
the planar reach but increases the vertical clearance required 
when the arm is folded into its home position.

The selection of these link lengths reflects a compromise between 
reach and load. Shorter links reduce the lever-arm torque on the 
shoulder joint and improve stability under load, while longer 
links extend the workspace radius. The chosen configuration 
provides a practical workspace radius of approximately 40~cm, 
which is sufficient for the workspace size used in the project 
(50~$\times$~30~cm), while keeping the load on the shoulder 
motor within the torque envelope of the XM540-W270 selected in 
the previous paragraph.

\hmsubsubsection{Embedded Motor Controller Selection}{UMA}{}

The motor controller is the embedded board responsible for 
translating commands received from the host processor into 
electrical signals on the Dynamixel TTL bus. Three candidate 
boards were evaluated during the early phase of the project: an 
Arduino Mega with a Dynamixel Shield, the ROBOTIS OpenCM~9.04, 
and the ROBOTIS OpenRB-150. The relevant specifications of the 
three candidates are summarised in Table~\ref{tab:controller_comparison}. 
The OpenRB-150 was selected as the final controller. The reasoning 
behind this choice is described below.

\begin{table}[H]
    \centering
    \renewcommand{\arraystretch}{1.3}
    \begin{tabular}{|p{3.3cm}|p{3.3cm}|p{3.3cm}|p{3.3cm}|}
        \hline
        \textbf{Specification} & \textbf{Arduino Mega + Shield} & 
        \textbf{OpenCM 9.04} & \textbf{OpenRB-150} \\
        \hline
        Processor & ATmega2560 (8-bit AVR) & STM32F103 (ARM Cortex-M3) & 
        SAMD21 (ARM Cortex-M0+) \\
        \hline
        Clock speed & 16 MHz & 72 MHz & 48 MHz \\
        \hline
        Architecture & 8-bit & 32-bit & 32-bit \\
        \hline
        Built-in TTL ports & 1 & 4 & 4 \\
        \hline
        Motor power input & External wiring & 5--16 V dedicated & 
        5--12.6 V dedicated \\
        \hline
        Host connector & USB-B & Micro-USB & USB-C \\
        \hline
        Arduino IDE support & Native & Native & Native \\
        \hline
    \end{tabular}
    \caption{Comparison of the three embedded motor controllers 
    evaluated for the project. Specifications taken from the 
    manufacturer documentation \cite{robotis_dynamixel}.}
    \label{tab:controller_comparison}
\end{table}

\hmsubsubsection{Initial Evaluation: Arduino Mega with Dynamixel Shield}{UMA}{}
The previous iteration of the project used an Arduino Mega 
combined with a Dynamixel Shield as the motor controller 
\cite{hivemind}. This combination was reviewed at the start of 
the project as a baseline, but several limitations were 
identified. The Arduino Mega is based on an 8-bit ATmega2560 
processor running at 16~MHz, which is significantly slower than 
the 32-bit microcontrollers used in the alternative boards. 
Lower processing speed translates into longer command-response 
cycles on the motor bus, which can introduce visible latency in 
the arm's motion when many commands are issued in sequence. In 
addition, the Dynamixel Shield does not provide a dedicated 
power input for the motor bus, requiring external power injection 
through additional wiring, and the previous iteration reported 
that the on-board regulator could not supply sufficient current 
for the analogue gripper servo, requiring an unstable 
breadboard-based workaround. Replacing the analogue gripper 
servo with a Dynamixel motor (Section~7.2.7) removed the immediate 
need for the PCA9685 driver, but the underlying limitations of 
the Arduino Mega remained.

\hmsubsubsection{Second Evaluation: OpenCM 9.04}{UMA}{}
The OpenCM~9.04, a microcontroller board developed by ROBOTIS 
specifically for Dynamixel control, was tested next. The board 
provides direct support for the Dynamixel protocol through the 
manufacturer's own libraries and offers four built-in TTL ports, 
addressing the wiring limitations of the Arduino Mega. However, 
two practical issues led to its rejection. First, the board's 
motor power input is designed for a battery connector rather 
than a regulated bench supply, which would have required an 
additional adapter to integrate with the laboratory power supply 
selected in Section~7.2.4. Second, the specific board acquired 
for the project ceased to respond to the host after the pin 
headers were soldered for connection to the motor bus, and the 
fault could not be recovered within the time budget allocated 
for controller selection. The OpenCM~9.04 was therefore set 
aside in favour of an alternative from the same manufacturer.

\hmsubsubsection{Selected Controller: OpenRB-150}{UMA}{}
The ROBOTIS OpenRB-150 was selected as the final motor controller. 
It is a more recent board developed by ROBOTIS as a successor to 
the OpenCM~9.04, and it shares the four built-in TTL ports while 
providing a motor power input compatible with the regulated 
laboratory bench supply, removing the need for the battery 
adapter that the OpenCM~9.04 would have required. The board 
introduces two further improvements that were directly relevant 
to the project.

The first improvement is the use of a USB-C connector for 
communication with the host processor, in place of the micro-USB 
connector used on the OpenCM~9.04. This was relevant for two 
reasons. The USB-C connector is mechanically more robust under 
repeated insertion and provides higher data throughput. In 
addition, the development environment used during the project 
was based on a macOS host, on which native USB-C ports are 
standard and which therefore required no adapter for board 
flashing or serial communication. This simplified the firmware 
development cycle substantially compared with the micro-USB 
alternatives.

The second improvement is the use of the SAMD21 processor, a 
32-bit ARM Cortex-M0+ running at 48~MHz. While this is nominally 
slower than the 72~MHz STM32F103 in the OpenCM~9.04, the SAMD21 
is well matched to the requirements of the application, providing 
sufficient throughput for sub-millisecond command-response cycles 
on the motor bus. Compared with the 8-bit ATmega2560 on the 
Arduino Mega, the SAMD21 produces noticeably smoother and more 
responsive arm motion, since the controller can process incoming 
position commands and issue corresponding TTL packets to the 
motor bus more rapidly.

The OpenRB-150 was tested immediately upon arrival and 
established reliable communication with all five Dynamixel 
servos on the first attempt. Based on this result, the board was 
adopted as the final motor controller and the Arduino Mega with 
Dynamixel Shield was retired from the design. The firmware 
running on the OpenRB-150 is described in 
Section~8.4, and the communication 
protocol between the host processor and the OpenRB-150 is 
described in Section~8.5.